\documentclass[12pt]{article}
\usepackage[utf8]{inputenc}
\usepackage[T1]{fontenc}
\usepackage[margin=1in]{geometry}
\usepackage{setspace}
\usepackage{xcolor}
\usepackage[colorlinks=true, urlcolor=blue!60!black, linkcolor=blue!60!black]{hyperref}
\setlength{\parindent}{0em}
\setlength{\parskip}{0.45em}
\renewcommand{\baselinestretch}{1.0}
\begin{document}
\begin{center}
{\Large \textit{CFPB v. Community Financial Services Ass'n of America, Ltd.}, 601 U.S. 416 (2024)\par}
\vspace{0.3em}
{\normalsize Case Brief --- One-Page Edition\par}
\end{center}

\textbf{Memory Jogger.} The CFPB's standing, capped draw on Federal Reserve earnings is a valid appropriation (\href{https://supreme.justia.com/cases/federal/us/601/22-448/}{Justia}).

\textbf{Facts.} Dodd-Frank lets the CFPB Director draw each year from the Federal Reserve's combined earnings the amount ``reasonably necessary'' to run the Bureau, up to an inflation-adjusted cap. Payday lenders challenged the Bureau's 2017 Payday Lending Rule.

\textbf{Procedural History.} The Fifth Circuit held the funding scheme violated the Appropriations Clause and vacated the rule.

\textbf{Votes.} Reversed 7--2: Thomas for the Court; Kagan and Jackson concurred; Alito, with Gorsuch, dissented.

\textbf{Holding.} An appropriation need only identify a source of public funds and authorize their expenditure for designated purposes.

\textbf{Reasoning.} Founding-era usage, the Army Clause's two-year limit, Parliament's capped grants, and the First Congress's fee-funded Customs Service and Post Office show no annual or fixed-sum requirement; Congress set the cap. Kagan stressed two centuries of practice. Alito said the Clause had become ``a minor vestige.''

\textbf{Examples.} \emph{Same-Side:} a fee-funded registry. \emph{Opposite-Side:} a permanent Army appropriation. \emph{Fence-Sitter:} an uncapped, perpetual draw for an AI agency.

\textbf{Critique.} Critics call the test toothless; defenders note Congress can repeal standing appropriations.

\textbf{Key Quotation.} ``Under the Appropriations Clause, an appropriation is simply a law that authorizes expenditures from a specified source of public money for designated purposes.'' --- Justice Thomas, 601 U.S. at 424.

\end{document}
